\section{Erd\H{o}s Problem 625 --- ROUND-2}

\subsection*{1) ROUND-2 OBJECTIVE}

\textbf{Path pursued: (C) obstruction / correction-by-reduction.}

Round-1 recorded the classical asymptotics and the 2024 progress of Heckel and Steiner,
but it still left a qualitative gap:
why do anti-concentration statements about $\chi(G(n,1/2))$ control the random variable
$\chi(G)-\zeta(G)$, and what exact obstruction remains to a full solution for \emph{all} $n$?

In this round I prove a precise quantitative transfer lemma:
small gap $\chi-\zeta$ forces near-concentration of $\chi$ in a short interval.
This isolates the remaining obstruction sharply.

\subsection*{2) Round-1 FOUNDATION USED}

I use the following Round-1 facts as fixed background:

\begin{itemize}
\item the definitions of $\chi(G)$ and $\zeta(G)$;
\item the classical bounds
\[
\frac{n}{2\log_2 n}\le \zeta(G)\le \chi(G)\le (1+o(1))\frac{n}{2\log_2 n}
\]
for $G\sim G(n,1/2)$;
\item the Round-1 summary of recent progress by Steiner and Heckel showing that
$\chi(G)-\zeta(G)$ is not bounded with high probability and that Heckel proves a positive answer
for roughly $95\%$ of all $n$.
\end{itemize}

\subsection*{3) NEW INSIGHT / TOOL (ROUND-2)}

The new tool is a \emph{quantitative concentration transfer lemma}.

For fixed $m\ge 0$, define
\[
q_{n,m}:=\Pr\bigl(\chi(G)-\zeta(G)>m\bigr),
\qquad G\sim G(n,1/2).
\]
I prove:

\begin{enumerate}
\item for every integer $a$,
\[
q_{n,m}\ge \Pr(\chi(G)\le a)\Pr(\chi(G)\ge a+m+1);
\]
\item consequently, there exists an interval of width $m$ that contains
$\chi(G)$ with probability at least $1-2\sqrt{q_{n,m}}$.
\end{enumerate}

So if $\chi(G)-\zeta(G)\le m$ with high probability, then $\chi(G)$ must be concentrated,
with high probability, in an interval of width $m$.
This is the exact obstruction that remains.

\subsection*{4) ATTACK PLAN (ROUND-2)}

The Round-1 obstacle was qualitative:
recent work implies large gaps for infinitely many $n$ and for about $95\%$ of all $n$,
but not for \emph{every} $n$.

To sharpen this, I prove a deterministic-probabilistic lemma reducing the open problem to one precise
question about fixed-width concentration of $\chi(G(n,1/2))$.
The proof uses only three facts:

\begin{enumerate}
\item $\zeta(G)\le \chi(\overline G)$;
\item $\overline G$ has the same distribution as $G$ when $G\sim G(n,1/2)$;
\item Harris--FKG for increasing events in the edge-exposure product space.
\end{enumerate}

\subsection*{5) WORK (ROUND-2)}

Let $G\sim G(n,1/2)$ and let $\overline G$ denote its complement.

\paragraph{Lemma 1 (deterministic comparison).}
For every graph $G$,
\[
\zeta(G)\le \chi(\overline G).
\]

\paragraph{Proof.}
A proper colouring of $\overline G$ partitions $V(G)$ into independent sets of $\overline G$,
equivalently into cliques of $G$.
Such a partition is a valid cocolouring of $G$.
Hence the minimum number of colours in a cocolouring is at most $\chi(\overline G)$. \hfill $\square$

\paragraph{Lemma 2 (quantitative transfer).}
Fix an integer $m\ge 0$ and put
\[
q_{n,m}:=\Pr\bigl(\chi(G)-\zeta(G)>m\bigr).
\]
Then for every integer $a$,
\[
q_{n,m}\ge \Pr(\chi(G)\le a)\Pr(\chi(G)\ge a+m+1).
\tag{1}
\]

\paragraph{Proof.}
Consider the events
\[
A_a:=\{\chi(G)\ge a+m+1\},
\qquad
B_a:=\{\chi(\overline G)\le a\}.
\]
The event $A_a$ is increasing in the edge set of $G$:
adding edges cannot decrease $\chi(G)$.

The event $B_a$ is also increasing in the edge set of $G$:
adding an edge to $G$ removes an edge from $\overline G$,
and removing edges cannot increase $\chi(\overline G)$.
Therefore $A_a$ and $B_a$ are both increasing events in the product probability space
of edge indicators of $G$.
By Harris--FKG,
\[
\Pr(A_a\cap B_a)\ge \Pr(A_a)\Pr(B_a).
\tag{2}
\]
Because $\overline G$ has the same distribution as $G(n,1/2)$,
\[
\Pr(B_a)=\Pr(\chi(G)\le a).
\tag{3}
\]
On the event $A_a\cap B_a$, Lemma 1 gives
\[
\zeta(G)\le \chi(\overline G)\le a,
\]
while also $\chi(G)\ge a+m+1$.
Hence
\[
\chi(G)-\zeta(G)\ge m+1,
\]
so $A_a\cap B_a\subseteq \{\chi(G)-\zeta(G)>m\}$.
Thus, using (2) and (3),
\[
q_{n,m}\ge \Pr(A_a\cap B_a)\ge \Pr(\chi(G)\ge a+m+1)\Pr(\chi(G)\le a),
\]
which is exactly (1). \hfill $\square$

\paragraph{Lemma 3 (concentration consequence).}
For every $n$ and every integer $m\ge 0$, there exists an integer $a_n$ such that
\[
\Pr\bigl(a_n\le \chi(G)\le a_n+m\bigr)\ge 1-2\sqrt{q_{n,m}}.
\tag{4}
\]

\paragraph{Proof.}
Let
\[
F(t):=\Pr(\chi(G)\le t).
\]
If $q_{n,m}=0$, choose $a_n$ minimal such that $F(a_n)>0$.
Then $F(a_n-1)=0$, and applying Lemma 2 with $a=a_n$ gives
$F(a_n)\bigl(1-F(a_n+m)\bigr)=0$.
Since $F(a_n)>0$, we get $F(a_n+m)=1$, so
\[
\Pr\bigl(a_n\le \chi(G)\le a_n+m\bigr)=1,
\]
which is exactly (4). Thus it remains only to treat the case $q_{n,m}>0$.

Choose $a_n$ minimal such that
\[
F(a_n)\ge \sqrt{q_{n,m}}.
\]
Then by minimality,
\[
F(a_n-1)<\sqrt{q_{n,m}}.
\tag{5}
\]
Apply Lemma 2 with $a=a_n$:
\[
q_{n,m}\ge F(a_n)\bigl(1-F(a_n+m)\bigr).
\]
Since $F(a_n)\ge \sqrt{q_{n,m}}$, it follows that
\[
1-F(a_n+m)\le \sqrt{q_{n,m}}.
\tag{6}
\]
Now
\[
\Pr\bigl(a_n\le \chi(G)\le a_n+m\bigr)
=
F(a_n+m)-F(a_n-1).
\]
Using (5) and (6),
\[
F(a_n+m)-F(a_n-1)\ge 1-\sqrt{q_{n,m}}-\sqrt{q_{n,m}}
=1-2\sqrt{q_{n,m}}.
\]
This proves (4). \hfill $\square$

\paragraph{Corollary 4.}
If for some fixed $m$ one has
\[
\Pr\bigl(\chi(G)-\zeta(G)\le m\bigr)\to 1,
\]
then there exists a sequence of intervals $I_n$ of width $m$ such that
\[
\Pr\bigl(\chi(G)\in I_n\bigr)\to 1.
\]

\paragraph{Proof.}
This is immediate from Lemma 3, since in that case $q_{n,m}\to 0$. \hfill $\square$

\paragraph{Corollary 5.}
If for some fixed $m$ and some $\delta>0$ every interval $I$ of width $m$ satisfies
\[
\Pr(\chi(G)\in I)\le 1-\delta,
\]
then
\[
\Pr\bigl(\chi(G)-\zeta(G)>m\bigr)\ge \delta^2/4.
\]

\paragraph{Proof.}
If $q_{n,m}<\delta^2/4$, then Lemma 3 would produce an interval of width $m$
containing $\chi(G)$ with probability $>1-\delta$, contradiction. \hfill $\square$

\paragraph{Interpretation.}
The full Erd\H{o}s--Gimbel problem is therefore reduced to the following sharply formulated obstruction:

\begin{quote}
\emph{For each fixed $m$, can $\chi(G(n,1/2))$ concentrate with probability $1-o(1)$
inside some interval of width $m$?}
\end{quote}

If the answer is \emph{no} for every fixed $m$, then $\chi(G)-\zeta(G)\to\infty$ with high probability.

\subsection*{6) ADVERSARIAL VERIFICATION}

\begin{itemize}
\item \emph{Monotonicity of $B_a=\{\chi(\overline G)\le a\}$.}
This is easy to get wrong because the event is phrased in terms of $\overline G$.
But as edges are added to $G$, edges are removed from $\overline G$, and removing edges cannot increase
chromatic number. So $B_a$ is indeed increasing in the coordinates of $G$.

\item \emph{Use of complement symmetry.}
For $G\sim G(n,1/2)$, the complement $\overline G$ has the same distribution as $G$.
This is essential in passing from $\Pr(B_a)$ to $\Pr(\chi(G)\le a)$.

\item \emph{Endpoint conventions.}
I formulated the gap event as $\chi(G)-\zeta(G)>m$ and the interval as $[a_n,a_n+m]$.
This matches the discrete inequality in Lemma 2:
the upper tail is $\Pr(\chi\ge a+m+1)$.

\item \emph{What the lemma does \emph{not} prove.}
Lemma 3 does not itself show that $\chi(G)-\zeta(G)\to\infty$ with high probability.
It shows instead that any failure of the conjecture must be accompanied by near-total concentration
of $\chi(G)$ in a bounded-width interval.
Thus it identifies the exact remaining obstruction.
\end{itemize}

\subsection*{7) FINAL}

\textbf{UNRESOLVED (BUT STRICTLY ADVANCED).}

The strict advance over Round-1 is a quantitative theorem:
\emph{small cochromatic gap forces short-interval concentration of the chromatic number.}
This turns the open problem into a sharp fixed-width concentration problem for $\chi(G(n,1/2))$.
That is a precise obstruction, rather than the looser heuristic connection recorded in Round-1.

\subsection*{8) COMPLETION ESTIMATE (MANDATORY)}

COMPLETION: 55\%

\subsection*{9) REFERENCES}

\begin{itemize}
\item T. F. Bloom, \emph{Erd\H{o}s Problem \#625}, \texttt{erdosproblems.com/625}, accessed 2026-03-07.
\item R. Steiner, \emph{On the difference between the chromatic and cochromatic number}, arXiv:2408.02400v2.
\item A. Heckel, \emph{On a question of Erd\H{o}s and Gimbel on the cochromatic number}, Electron. J. Combin. \textbf{31} (2024), P4.72.
\item A. Heckel, \emph{The difference between the chromatic and the cochromatic number of a random graph}, arXiv:2409.17614v2.
\item N. Alon and J. H. Spencer, \emph{The Probabilistic Method}, for the Harris--FKG inequality.
\end{itemize}

\subsection*{10) OUTPUT FORMAT}

This section is written in drop-in \LaTeX\ form and starts directly from \verb|\section{}|, with no preamble.
